\documentclass[10pt,landscape,a4paper]{article}
\usepackage[top=0.3cm,bottom=0.3cm,left=0.3cm,right=0.3cm]{geometry}
\usepackage{multicol}
\usepackage{amsmath,amssymb}
\usepackage{array}
\usepackage[dvipsnames]{xcolor}
\setlength{\columnseprule}{0.4pt}
\setlength{\columnsep}{0.7em}
\newcommand{\problem}[1]{\textbf{\colorbox{cyan!15}{#1}}}
\newcommand{\answer}[1]{\colorbox{green!15}{#1}}
\pagestyle{empty}
\parindent=0pt

\begin{document}
\begin{multicols}{3}
\scriptsize

\section*{Worked Examples}

\subsection*{Complex Arithmetic}

\problem{Ex 1: Simplify $(1-i)^{12}$}

\textbf{Step 1:} Convert to polar
\begin{itemize}
\item $|1-i| = \sqrt{1^2 + 1^2} = \sqrt{2}$
\item Angle: $\theta = -\pi/4$ (4th quadrant)
\item So $1-i = \sqrt{2}e^{-i\pi/4}$
\end{itemize}

\textbf{Step 2:} Apply power rule
\begin{itemize}
\item $(1-i)^{12} = (\sqrt{2})^{12} \cdot e^{-12i\pi/4}$
\item $= 2^6 \cdot e^{-3i\pi} = 64 \cdot e^{-3i\pi}$
\end{itemize}

\textbf{Step 3:} Convert back
\begin{itemize}
\item $e^{-3i\pi} = e^{-i\pi} = -1$
\item Final: $64 \cdot (-1) = \answer{-64}$
\end{itemize}

\problem{Ex 2: Complex product $(1/2+\sqrt{3}/2\cdot i)(1/2-\sqrt{3}/2\cdot i)$}

\textbf{Method 1 - Conjugate Pattern:}
\begin{itemize}
\item Pattern: $(a+bi)(a-bi) = a^2 + b^2$
\item Here: $a = 1/2$, $b = \sqrt{3}/2$
\item Result: $(1/2)^2 + (\sqrt{3}/2)^2 = 1/4 + 3/4 = \answer{1}$
\end{itemize}

\textbf{Method 2 - FOIL (verify):}
\begin{itemize}
\item Real part: $(1/2)(1/2) - (\sqrt{3}/2)(-\sqrt{3}/2)$
\item $= 1/4 + 3/4 = 1$
\item Imaginary: $(1/2)(-\sqrt{3}/2) + (\sqrt{3}/2)(1/2) = 0$
\item Result: $1 + 0i = \answer{1}$
\end{itemize}

\problem{Ex 3: Rotate $-1/2+\sqrt{3}/2\cdot i$ by $-\pi/6$}

\textbf{Step 1:} Identify original angle
\begin{itemize}
\item This is $e^{i2\pi/3}$ (120 degrees)
\item Verify: $\cos(2\pi/3) = -1/2$, $\sin(2\pi/3) = \sqrt{3}/2$
\end{itemize}

\textbf{Step 2:} Apply rotation
\begin{itemize}
\item Multiply by $e^{-i\pi/6}$ to rotate clockwise
\item $e^{i2\pi/3} \cdot e^{-i\pi/6} = e^{i(2\pi/3 - \pi/6)}$
\item $= e^{i(4\pi/6 - \pi/6)} = e^{i\pi/2}$
\end{itemize}

\textbf{Step 3:} Convert to rectangular
\begin{itemize}
\item $e^{i\pi/2} = \cos(\pi/2) + i\sin(\pi/2) = 0 + i = \answer{i}$
\end{itemize}

\subsection*{Polynomial Factoring}

\problem{Factor $x^3 + 1$ and find all roots}

\textbf{Step 1:} Use sum of cubes formula
\begin{itemize}
\item $x^3 + 1 = (x+1)(x^2-x+1)$
\item First root: $x = -1$
\end{itemize}

\textbf{Step 2:} Quadratic formula on $x^2 - x + 1 = 0$
\begin{itemize}
\item $x = \frac{1 \pm \sqrt{1-4}}{2} = \frac{1 \pm \sqrt{-3}}{2}$
\item $= \frac{1 \pm i\sqrt{3}}{2}$
\end{itemize}

\textbf{Step 3:} All three roots
\begin{itemize}
\item Root 1: $x = -1$
\item Root 2: $x = \frac{1 + i\sqrt{3}}{2} = e^{i\pi/3}$
\item Root 3: $x = \frac{1 - i\sqrt{3}}{2} = e^{-i\pi/3}$
\item These are the three cube roots of -1
\end{itemize}

\subsection*{Aliasing}

\problem{Find alias: $f_s = 44100$ Hz, signal = 32100 Hz}

\textbf{Step 1:} Check Nyquist
\begin{itemize}
\item Nyquist freq = $f_s/2 = 22050$ Hz
\item 32100 $>$ 22050, so aliasing occurs
\end{itemize}

\textbf{Step 2:} Calculate alias frequency
\begin{itemize}
\item Use: $f_{alias} = |f - nf_s|$ where $n$ minimizes
\item $n=0$: $|32100 - 0| = 32100$ Hz
\item $n=1$: $|32100 - 44100| = 12000$ Hz (smaller!)
\item $n=2$: $|32100 - 88200| = 56100$ Hz (larger)
\end{itemize}

\textbf{Step 3:} Verify and conclude
\begin{itemize}
\item 12000 Hz is in valid range [0, 22050]
\item The 32100 Hz signal appears as \answer{12000 Hz}
\end{itemize}

\subsection*{DFT Calculation}

\problem{Find $X[1]$ for $x = [1, 0, 1]$, $N=3$}

\textbf{Step 1:} Write DFT formula
$$X[k] = \sum_{n=0}^{N-1} x[n]e^{-i2\pi kn/N}$$

\textbf{Step 2:} Substitute values ($k=1$, $N=3$)
\begin{itemize}
\item $X[1] = x[0]e^0 + x[1]e^{-i2\pi/3} + x[2]e^{-i4\pi/3}$
\item $= 1(1) + 0(e^{-i2\pi/3}) + 1(e^{-i4\pi/3})$
\item $= 1 + e^{-i4\pi/3}$
\end{itemize}

\textbf{Step 3:} Simplify the exponential
\begin{itemize}
\item $e^{-i4\pi/3} = \cos(-4\pi/3) + i\sin(-4\pi/3)$
\item $= \cos(4\pi/3) - i\sin(4\pi/3)$
\item $= -1/2 - i\cdot(-\sqrt{3}/2) = -1/2 + i\sqrt{3}/2$
\end{itemize}

\textbf{Step 4:} Final answer
\begin{itemize}
\item $X[1] = 1 + (-1/2 + i\sqrt{3}/2) = 1/2 + i\sqrt{3}/2$
\item This equals $e^{i\pi/3}$
\item DFT coefficient: $c_1 = X[1]/N = \answer{e^{i\pi/3}/3}$
\end{itemize}

\subsection*{FIR Filter Analysis}

\problem{Analyze $y_t = x_t - (1/4)x_{t-1}$}

\textbf{Step 1:} Transfer function
\begin{itemize}
\item Direct: $H(z) = 1 - (1/4)z^{-1}$
\end{itemize}

\textbf{Step 2:} Find zeros and poles
\begin{itemize}
\item Zero: Set $1 - (1/4)z^{-1} = 0$
\item Solve: $z^{-1} = 4$, so $z = 1/4$
\item Pole: At $z = 0$ (from $z^{-1}$ term)
\item \answer{Zero at $z=1/4$, pole at $z=0$}
\end{itemize}

\textbf{Step 3:} Frequency response at $\omega = 0$
\begin{itemize}
\item $H(0) = 1 - (1/4)e^{-i \cdot 0} = 1 - 1/4 = \answer{3/4}$
\item DC gain is 0.75 (attenuation)
\end{itemize}

\textbf{Step 4:} Magnitude at $\omega = \pi/2$
\begin{itemize}
\item $H(\pi/2) = 1 - (1/4)e^{-i\pi/2}$
\item $e^{-i\pi/2} = -i$, so $H(\pi/2) = 1 - (1/4)(-i) = 1 + i/4$
\item $|H(\pi/2)| = |1 + i/4| = \sqrt{1 + 1/16} = \answer{\sqrt{17}/4}$
\end{itemize}

\subsection*{IIR Filter with Impulse}

\problem{Analyze $y_t = x_t - (1/4)y_{t-1}$, impulse response}

\textbf{Step 1:} Find transfer function
\begin{itemize}
\item Rearrange: $y_t + (1/4)y_{t-1} = x_t$
\item Z-transform: $Y(z)[1 + (1/4)z^{-1}] = X(z)$
\item $H(z) = 1/[1 + (1/4)z^{-1}] = z/(z + 1/4)$
\end{itemize}

\textbf{Step 2:} Check stability
\begin{itemize}
\item Pole at $z = -1/4$
\item $|pole| = |-1/4| = 1/4 < 1$ \answer{STABLE}
\end{itemize}

\textbf{Step 3:} Impulse response calculation
\begin{itemize}
\item Input: $x = [1, 0, 0, ...]$
\item $y_0 = 1 - (1/4)(0) = 1$
\item $y_1 = 0 - (1/4)(1) = -1/4$
\item $y_2 = 0 - (1/4)(-1/4) = 1/16$
\item $y_3 = 0 - (1/4)(1/16) = -1/64$
\end{itemize}

\textbf{Step 4:} Pattern recognition
\begin{itemize}
\item $h[n] = (-1/4)^n$ for $n \geq 0$
\item Alternating sign, exponential decay
\item Sum: $\sum h[n] = 1/(1-(-1/4)) = \answer{4/5}$
\end{itemize}

\subsection*{Reson Filter Design}

\problem{Design reson: $\omega = \pi/4$, bandwidth $B = 0.1$}

\textbf{Step 1:} Calculate pole radius from bandwidth
\begin{itemize}
\item Narrow band formula: $B \approx 2(1-R)$
\item $0.1 = 2(1-R)$, so $1-R = 0.05$
\item $R = 0.95$
\end{itemize}

\textbf{Step 2:} Place complex conjugate poles
\begin{itemize}
\item Poles at $z = 0.95e^{\pm i\pi/4}$
\item In rectangular: $0.95 \cdot (\frac{1}{\sqrt{2}} \pm i\frac{1}{\sqrt{2}})$
\item $= \frac{0.95}{\sqrt{2}}(1 \pm i) = 0.672(1 \pm i)$
\end{itemize}

\textbf{Step 3:} Write transfer function
\begin{itemize}
\item $H(z) = \frac{1}{1 - 2R\cos(\pi/4)z^{-1} + R^2z^{-2}}$
\item $\cos(\pi/4) = 1/\sqrt{2}$, so $2R\cos(\pi/4) = 2(0.95)/\sqrt{2} = 1.343$
\item $H(z) = \frac{1}{1 - 1.343z^{-1} + 0.9025z^{-2}}$
\end{itemize}

\textbf{Step 4:} Calculate peak gain
\begin{itemize}
\item At resonance ($\omega = \pi/4$): 
\item Approximate gain $\approx 1/(1-R) = 1/0.05 = 20$
\item In dB: $20\log_{10}(20) = \answer{26 dB}$
\end{itemize}

\subsection*{Plucked String Tuning}

\problem{Tune to 440 Hz, $f_s = 44100$ Hz}

\textbf{Step 1:} Calculate ideal period in samples
\begin{itemize}
\item $P = f_s/f_0 = 44100/440 = 100.227$ samples
\item Need to achieve this total delay
\end{itemize}

\textbf{Step 2:} Allocate delays
\begin{itemize}
\item Try integer delay $L = 100$
\item Low-pass filter adds 0.5 samples
\item Total so far: $100 + 0.5 = 100.5$
\item Still need: $100.227 - 100.5 = -0.273$ (negative!)
\end{itemize}

\textbf{Step 3:} Adjust allocation
\begin{itemize}
\item Use $L = 99$ instead
\item With LP: $99 + 0.5 = 99.5$
\item Need from all-pass: $100.227 - 99.5 = 0.727$
\end{itemize}

\textbf{Step 4:} Calculate all-pass coefficient
\begin{itemize}
\item For delay $\delta = 0.727$:
\item $a = (1-\delta)/(1+\delta) = 0.273/1.727 = \answer{0.158}$
\end{itemize}

\textbf{Step 5:} Verify total delay
\begin{itemize}
\item Integer: 99 samples
\item Low-pass: 0.5 samples
\item All-pass: 0.727 samples
\item Total: $99 + 0.5 + 0.727 = 100.227$ \answer{Perfect!}
\end{itemize}

\subsection*{All-Pass Identification}

\problem{Is $H(z) = (0.5z + 1)/(z + 0.5)$ all-pass?}

\textbf{Step 1:} Convert to standard form
\begin{itemize}
\item Divide num and denom by $z$:
\item $H(z) = (0.5 + z^{-1})/(1 + 0.5z^{-1})$
\end{itemize}

\textbf{Step 2:} Check all-pass structure
\begin{itemize}
\item All-pass form: $H(z) = (a + z^{-1})/(1 + az^{-1})$
\item Our filter: numerator has $a = 0.5$
\item Denominator also has $a = 0.5$
\item Structure matches! \answer{YES, all-pass}
\end{itemize}

\textbf{Step 3:} Verify magnitude = 1
\begin{itemize}
\item At any frequency: $|H(e^{i\omega})| = |(0.5 + e^{-i\omega})/(1 + 0.5e^{-i\omega})|$
\item Numerator and denominator have same structure
\item Can show $|H(e^{i\omega})| = 1$ for all $\omega$
\end{itemize}

\textbf{Step 4:} Phase delay property
\begin{itemize}
\item All-pass changes phase without changing magnitude
\item Used for fractional delay adjustment
\item Phase: $\phi(\omega) = -2\arctan(\frac{0.5\sin\omega}{1+0.5\cos\omega})$
\end{itemize}

\end{multicols}
\end{document}